\section{Introduction}
Online shopping relies on images to convey shape, color, material, and fit, which creates a barrier for Blind and Low-Vision (BLV) users who browse with screen readers. Alternative text (alt text) is ``descriptive text that conveys the meaning of an image in digital content'' \cite{whatalttext}; it is the main channel through which a BLV shopper can access what a product photograph shows.

In practice, e-commerce alt text is frequently insufficient. Many sites duplicate the product title, paste fragments of the description, or, at best, run a generic image captioner. Because large marketplaces let third-party sellers supply their own metadata, alt-text quality also varies widely between listings. This matters because BLV shoppers specifically want appearance attributes such as color, logo, shape, and size \cite{wang2021}, and generic alt text does not support the decision to buy or not.

Vision Language Models (VLMs) can now describe images fluently, but good alt text depends on the image's context and purpose \cite{nngroup2024}: the same image may need a completely different description depending on where it appears and what the user is trying to do \cite{stangl2021}. W3C guidance makes the same point and adds that if nearby text already conveys the visual information, repeating it in the text alternative is redundant \cite{w3cimages,w3cinformative}. On a product page, the title and attribute list already tell the user that an item is, say, a black women's sweater. The useful question is therefore not ``what is in this image?'' but ``what visual information is still missing from this page?''

We introduce \sys (Visual Gap-Aware ALT Text Generation), a multi-stage pipeline that answers that question. \sys extracts the product page's Document Object Model (DOM) into structured context with a deterministic parser; all subsequent generation, validation, and rewriting steps are performed by Large Language Models (LLMs). Our contributions are:
\begin{itemize}
\item a gap-aware generation pipeline that decomposes a candidate description into atomic claims and checks them with three parallel validators (redundancy, accuracy, completeness) before a constrained fuser writes the final alt text;
\item a paired, claim-level comparison against a reproduced context-aware baseline \cite{mohanbabu2024} under a controlled model backend, with Wilcoxon signed-rank tests and Holm correction;
\item a fusion-stage ablation that isolates what each validator contributes to the final text.
\end{itemize}
